\chapter*{PERNYATAAN PENGGUNAAN AI}
\addcontentsline{toc}{chapter}{PERNYATAAN PENGGUNAAN AI}
Saya yang bertanda tangan di bawah ini menyatakan bahwa dalam penulisan Proposal Tugas Akhir ini, saya telah menggunakan alat bantu kecerdasan artifisial (AI) untuk beberapa aspek tertentu. Berikut adalah rincian penggunaan AI dalam penulisan tugas akhir ini.
\begin{table}[h]
\centering
\small
\begin{tabularx}{\textwidth}{|p{0.5cm}| p{4cm} | p{1.8cm} | p{1.3cm} | X | }
\hline
\textbf{No.} & \textbf{Jenis} & \textbf{Alat Bantu} & \textbf{Tingkat} & \textbf{Tempat}\\
\hline
1 & Memeriksa ejaan dan tata bahasa & DeepL & Rendah & Bab I--III (pengecekan padanan istilah dari referensi berbahasa Inggris) \\
\hline
2 & Pembuatan grafik, gambar, atau ilustrasi & Gemini, Claude & Rendah & Bab III dan Bab IV (penyusunan sintaks diagram alur berbasis TikZ LaTeX) \\
\hline
3 & Pencarian informasi atau referensi & Paper Guide, Gemini & Sedang & Bab I dan Bab II (penelusuran dan pemetaan literatur terkait \textit{neuro-NLP} serta \textit{brain-tuning}) \\
\hline
4 & Menjelaskan definisi, konsep, gambar, tabel, atau topik tertentu & Paper Guide, Claude & Sedang & Bab II (memahami istilah neurosains, fMRI, dan metode pada artikel rujukan) \\
\hline
5 & Penyusunan kode program & Gemini, Claude & Rendah & Bab III, Bab IV, dan Bab V (penulisan kode format tabel LaTeX dan TikZ) \\
\hline
6 & Penyusunan formula atau perhitungan matematis & Claude & Rendah & Bab II dan Bab V (penulisan notasi LaTeX formula serta persamaan estimasi anggaran waktu GPU) \\
\hline
\end{tabularx}
\end{table}


Bandung, \today\\
\\[2cm]

\underline{Ananda  Fajrul Zidan} \\
NIM: 18223067